\documentclass{article}
\usepackage{lmodern}
\usepackage{amsmath}
\usepackage{listings}
\usepackage{fullpage}
\usepackage{parskip}
\usepackage{xcolor}
\lstset{
	basicstyle=\ttfamily,
	columns=fullflexible,
	frame=single,
	breaklines=true,
	postbreak=\mbox{\textcolor{red}{$\hookrightarrow$}\space},
}
\begin{document}
\title{Experiment `p\_4\_5\_analyse\_floats\_array' Results}
\author{\today}
\date{}
\maketitle
\textbf{Experiment outcome: }SUCCESS
\\\textbf{Bad responses: }0
\\\textbf{Responses containing}~\texttt{assume}~\textbf{: }0
\\\textbf{Responses containing}~\texttt{expect}~\textbf{: }0
\\\textbf{Resolution attempts: }3
\\\textbf{Hard fails (resolution): }2
\\\textbf{Soft fails (resolution): }0
\\\textbf{Verification attempts: }1
\section*{Problem Specification}
\textbf{Problem name: }p\_4\_5\_analyse\_floats\_array
\\\textbf{Natural language statement: }null
\\\textbf{Method signature: }p\_4\_5\_analyse\_floats\_array(values: array<real>) returns (average: real, min: real, max: real, range: real)
\subsection*{Ensures}
\begin{itemize}
\item \texttt{average == (sum(values[..]) / values.Length as real)}
\item \texttt{forall v :: v in values[..] ==> min <= v}
\item \texttt{forall v :: v in values[..] ==> max >= v}
\item \texttt{range == max - min}
\item \texttt{exists v :: v in values[..] \&\& min == v}
\item \texttt{exists v :: v in values[..] \&\& max == v}
\end{itemize}
\subsection*{Requires}
\begin{itemize}
\item \texttt{0 < values.Length}
\end{itemize}
\subsection*{Functional Code Given}
\begin{lstlisting}
function sum(values: seq<real>): real
{
  if |values| == 0 then 0.0
  else values[0] + sum(values[1..])
}
\end{lstlisting}
\clearpage
\section*{GenAI interactions}
Below you will find all interactions between the `user' (program) and the `assistant' (OpenAI).
\subsection*{Program $\rightarrow$ GenAI}
\begin{lstlisting}
You are given the following task to perform in Dafny:

Write a method to satisfy the below requirements.

The signature should be:

p_4_5_analyse_floats_array(values: array<real>) returns (average: real, min: real, max: real, range: real)

The method should respect the following contract:

requires 0 < values.Length, ensures average == (sum(values[..]) / values.Length as real), ensures forall v :: v in values[..] ==> min <= v, ensures forall v :: v in values[..] ==> max >= v, ensures range == max - min, ensures exists v :: v in values[..] && min == v, ensures exists v :: v in values[..] && max == v

The contract uses the following dafny code:

function sum(values: seq<real>): real
{
  if |values| == 0 then 0.0
  else values[0] + sum(values[1..])
}

These function(s) must not be used in your implementation of the method.

Produce and show only the Dafny body of this method, including the curly braces that surround it. Do not show the signature nor contract. You must not use 'assume' nor 'expect' anywhere in your code.
\end{lstlisting}
\subsection*{GenAI $\rightarrow$ Program}
\textbf{System fingerprint: }null
\\\textbf{ID: }chatcmpl-DnM9Bw34NxU9eIMyDa18zxnByCahV
\\\textbf{Model: }gpt-5.5-2026-04-23
\\\textbf{Created at: }1780655577
\\\textbf{Finish reason: }stop
\\\textbf{Usage: }558 tokens in, and 8302 tokens out
\begin{lstlisting}
{
  var n := values.Length;

  var total := 0.0;
  var i := n;
  while 0 < i
    invariant 0 <= i <= n
    invariant total == sum(values[i..n])
  {
    i := i - 1;
    assert 0 <= i < n;
    assert total == sum(values[i + 1..n]);
    assert |values[i..n]| > 0;
    assert (values[i..n])[0] == values[i];
    assert (values[i..n])[1..] == values[i + 1..n];
    assert sum(values[i..n]) == values[i] + sum(values[i + 1..n]);
    total := values[i] + total;
    assert total == sum(values[i..n]);
  }
  assert i == 0;
  assert values[0..n] == values[..];
  assert total == sum(values[..]);

  min := values[0];
  max := values[0];
  var minIndex := 0;
  var maxIndex := 0;
  i := 1;
  while i < n
    invariant 1 <= i <= n
    invariant forall k: int :: 0 <= k < i ==> min <= values[k]
    invariant forall k: int :: 0 <= k < i ==> max >= values[k]
    invariant 0 <= minIndex < i
    invariant min == values[minIndex]
    invariant 0 <= maxIndex < i
    invariant max == values[maxIndex]
  {
    var x := values[i];
    if x < min {
      min := x;
      minIndex := i;
    }
    if x > max {
      max := x;
      maxIndex := i;
    }
    i := i + 1;
  }
  assert i == n;
  assert forall k: int :: 0 <= k < n ==> min <= values[k];
  assert forall k: int :: 0 <= k < n ==> max >= values[k];

  average := total / (n as real);
  range := max - min;
  assert n == values.Length;
  assert average == sum(values[..]) / (values.Length as real);

  assert forall v: real :: v in values[..] ==> min <= v by {
    if v in values[..] {
      assert exists k: int :: 0 <= k < |values[..]| && (values[..])[k] == v;
      var k: int :| 0 <= k < |values[..]| && (values[..])[k] == v;
      assert |values[..]| == n;
      assert values[..][k] == values[k];
      assert v == values[k];
      assert min <= values[k];
    }
  }
  assert forall v: real :: v in values[..] ==> max >= v by {
    if v in values[..] {
      assert exists k: int :: 0 <= k < |values[..]| && (values[..])[k] == v;
      var k: int :| 0 <= k < |values[..]| && (values[..])[k] == v;
      assert |values[..]| == n;
      assert values[..][k] == values[k];
      assert v == values[k];
      assert max >= values[k];
    }
  }

  assert 0 <= minIndex < |values[..]|;
  assert values[..][minIndex] == min;
  assert min in values[..];
  assert exists v: real :: v in values[..] && min == v;

  assert 0 <= maxIndex < |values[..]|;
  assert values[..][maxIndex] == max;
  assert max in values[..];
  assert exists v: real :: v in values[..] && max == v;
}
\end{lstlisting}
\subsection*{Program $\rightarrow$ GenAI}
\begin{lstlisting}
When using dafny resolve, the below error is emitted and resolve fails:

problem_data/p_4_5_analyse_floats_array/program.dfy(73,7): Error: unresolved identifier: v
   |
73 |     if v in values[..] {
   |        ^

problem_data/p_4_5_analyse_floats_array/program.dfy(74,74): Error: unresolved identifier: v
   |
74 |       assert exists k: int :: 0 <= k < |values[..]| && (values[..])[k] == v;
   |                                                                           ^

problem_data/p_4_5_analyse_floats_array/program.dfy(75,64): Error: unresolved identifier: v
   |
75 |       var k: int :| 0 <= k < |values[..]| && (values[..])[k] == v;
   |                                                                 ^

problem_data/p_4_5_analyse_floats_array/program.dfy(78,13): Error: unresolved identifier: v
   |
78 |       assert v == values[k];
   |              ^

problem_data/p_4_5_analyse_floats_array/program.dfy(83,7): Error: unresolved identifier: v
   |
83 |     if v in values[..] {
   |        ^

problem_data/p_4_5_analyse_floats_array/program.dfy(84,74): Error: unresolved identifier: v
   |
84 |       assert exists k: int :: 0 <= k < |values[..]| && (values[..])[k] == v;
   |                                                                           ^

problem_data/p_4_5_analyse_floats_array/program.dfy(85,64): Error: unresolved identifier: v
   |
85 |       var k: int :| 0 <= k < |values[..]| && (values[..])[k] == v;
   |                                                                 ^

problem_data/p_4_5_analyse_floats_array/program.dfy(88,13): Error: unresolved identifier: v
   |
88 |       assert v == values[k];
   |              ^

8 resolution/type errors detected in program.dfy

Correct the error in code by altering only the method body. Produce and show only the Dafny body, including the curly braces that surround it. Do not show the signature nor contract. You must not use 'assume' nor 'expect' anywhere in your code.
\end{lstlisting}
\subsection*{GenAI $\rightarrow$ Program}
\textbf{System fingerprint: }null
\\\textbf{ID: }chatcmpl-DnMBykXOfm2DvAaqbWCCKlezP1kMg
\\\textbf{Model: }gpt-5.5-2026-04-23
\\\textbf{Created at: }1780655750
\\\textbf{Finish reason: }stop
\\\textbf{Usage: }2226 tokens in, and 1894 tokens out
\begin{lstlisting}
{
  var n := values.Length;

  var total := 0.0;
  var i := n;
  while 0 < i
    invariant 0 <= i <= n
    invariant total == sum(values[i..n])
  {
    i := i - 1;
    assert 0 <= i < n;
    assert total == sum(values[i + 1..n]);
    assert |values[i..n]| > 0;
    assert (values[i..n])[0] == values[i];
    assert (values[i..n])[1..] == values[i + 1..n];
    assert sum(values[i..n]) == values[i] + sum(values[i + 1..n]);
    total := values[i] + total;
    assert total == sum(values[i..n]);
  }
  assert i == 0;
  assert values[0..n] == values[..];
  assert total == sum(values[..]);

  min := values[0];
  max := values[0];
  var minIndex := 0;
  var maxIndex := 0;
  i := 1;
  while i < n
    invariant 1 <= i <= n
    invariant forall k: int :: 0 <= k < i ==> min <= values[k]
    invariant forall k: int :: 0 <= k < i ==> max >= values[k]
    invariant 0 <= minIndex < i
    invariant min == values[minIndex]
    invariant 0 <= maxIndex < i
    invariant max == values[maxIndex]
  {
    var x := values[i];
    if x < min {
      min := x;
      minIndex := i;
    }
    if x > max {
      max := x;
      maxIndex := i;
    }
    i := i + 1;
  }
  assert i == n;
  assert forall k: int :: 0 <= k < n ==> min <= values[k];
  assert forall k: int :: 0 <= k < n ==> max >= values[k];

  average := total / (n as real);
  range := max - min;
  assert n == values.Length;
  assert average == sum(values[..]) / (values.Length as real);

  assert forall v: real | v in values[..]
    ensures min <= v
  {
    assert exists k: int :: 0 <= k < |values[..]| && (values[..])[k] == v;
    var k: int :| 0 <= k < |values[..]| && (values[..])[k] == v;
    assert |values[..]| == n;
    assert (values[..])[k] == values[k];
    assert v == values[k];
    assert min <= values[k];
  }
  assert forall v: real | v in values[..]
    ensures max >= v
  {
    assert exists k: int :: 0 <= k < |values[..]| && (values[..])[k] == v;
    var k: int :| 0 <= k < |values[..]| && (values[..])[k] == v;
    assert |values[..]| == n;
    assert (values[..])[k] == values[k];
    assert v == values[k];
    assert max >= values[k];
  }

  assert 0 <= minIndex < |values[..]|;
  assert (values[..])[minIndex] == min;
  assert min in values[..];
  assert exists v: real :: v in values[..] && min == v;

  assert 0 <= maxIndex < |values[..]|;
  assert (values[..])[maxIndex] == max;
  assert max in values[..];
  assert exists v: real :: v in values[..] && max == v;
}
\end{lstlisting}
\subsection*{Program $\rightarrow$ GenAI}
\begin{lstlisting}
When using dafny resolve, the below error is emitted and resolve fails:

problem_data/p_4_5_analyse_floats_array/program.dfy(102,0): Error: invalid UnaryExpression
    |
102 | }
    | ^

problem_data/p_4_5_analyse_floats_array/program.dfy(101,54): Error: expected either 'by' or a semicolon following the assert expression
    |
101 |   assert exists v: real :: v in values[..] && max == v;
    |                                                       ^

problem_data/p_4_5_analyse_floats_array/program.dfy(101,54): Error: expected either 'by' or a semicolon following the assert expression
    |
101 |   assert exists v: real :: v in values[..] && max == v;
    |                                                       ^

3 parse errors detected in program.dfy

Correct the error in code by altering only the method body. Produce and show only the Dafny body, including the curly braces that surround it. Do not show the signature nor contract. You must not use 'assume' nor 'expect' anywhere in your code.
\end{lstlisting}
\subsection*{GenAI $\rightarrow$ Program}
\textbf{System fingerprint: }null
\\\textbf{ID: }chatcmpl-DnMCiRHOL2DrFiPseTRFDQk5jMVN1
\\\textbf{Model: }gpt-5.5-2026-04-23
\\\textbf{Created at: }1780655796
\\\textbf{Finish reason: }stop
\\\textbf{Usage: }3094 tokens in, and 2444 tokens out
\begin{lstlisting}
{
  var n := values.Length;

  var total := 0.0;
  var i := n;
  while 0 < i
    invariant 0 <= i <= n
    invariant total == sum(values[i..n])
  {
    i := i - 1;
    assert 0 <= i < n;
    assert total == sum(values[i + 1..n]);
    assert |values[i..n]| > 0;
    assert (values[i..n])[0] == values[i];
    assert (values[i..n])[1..] == values[i + 1..n];
    assert sum(values[i..n]) == values[i] + sum(values[i + 1..n]);
    total := values[i] + total;
    assert total == sum(values[i..n]);
  }
  assert i == 0;
  assert values[0..n] == values[..];
  assert total == sum(values[..]);

  min := values[0];
  max := values[0];
  var minIndex := 0;
  var maxIndex := 0;
  i := 1;
  while i < n
    invariant 1 <= i <= n
    invariant forall k: int :: 0 <= k < i ==> min <= values[k]
    invariant forall k: int :: 0 <= k < i ==> max >= values[k]
    invariant 0 <= minIndex < i
    invariant min == values[minIndex]
    invariant 0 <= maxIndex < i
    invariant max == values[maxIndex]
  {
    var x := values[i];
    if x < min {
      min := x;
      minIndex := i;
    }
    if x > max {
      max := x;
      maxIndex := i;
    }
    i := i + 1;
  }
  assert i == n;
  assert forall k: int :: 0 <= k < n ==> min <= values[k];
  assert forall k: int :: 0 <= k < n ==> max >= values[k];

  average := total / (n as real);
  range := max - min;
  assert n == values.Length;
  assert average == sum(values[..]) / (values.Length as real);

  forall v: real | v in values[..]
    ensures min <= v
  {
    assert exists k: int :: 0 <= k < |values[..]| && (values[..])[k] == v;
    var k: int :| 0 <= k < |values[..]| && (values[..])[k] == v;
    assert |values[..]| == n;
    assert (values[..])[k] == values[k];
    assert v == values[k];
    assert min <= values[k];
  }
  forall v: real | v in values[..]
    ensures max >= v
  {
    assert exists k: int :: 0 <= k < |values[..]| && (values[..])[k] == v;
    var k: int :| 0 <= k < |values[..]| && (values[..])[k] == v;
    assert |values[..]| == n;
    assert (values[..])[k] == values[k];
    assert v == values[k];
    assert max >= values[k];
  }

  assert 0 <= minIndex < |values[..]|;
  assert (values[..])[minIndex] == min;
  assert min in values[..];
  assert exists v: real :: (v in values[..]) && min == v;

  assert 0 <= maxIndex < |values[..]|;
  assert (values[..])[maxIndex] == max;
  assert max in values[..];
  assert exists v: real :: (v in values[..]) && max == v;
}
\end{lstlisting}
\section*{Final Program}
\begin{lstlisting}
function sum(values: seq<real>): real
{
  if |values| == 0 then 0.0
  else values[0] + sum(values[1..])
}

method p_4_5_analyse_floats_array(values: array<real>) returns (average: real, min: real, max: real, range: real)
	requires 0 < values.Length
	ensures average == (sum(values[..]) / values.Length as real)
	ensures forall v :: v in values[..] ==> min <= v
	ensures forall v :: v in values[..] ==> max >= v
	ensures range == max - min
	ensures exists v :: v in values[..] && min == v
	ensures exists v :: v in values[..] && max == v
{
  var n := values.Length;

  var total := 0.0;
  var i := n;
  while 0 < i
    invariant 0 <= i <= n
    invariant total == sum(values[i..n])
  {
    i := i - 1;
    assert 0 <= i < n;
    assert total == sum(values[i + 1..n]);
    assert |values[i..n]| > 0;
    assert (values[i..n])[0] == values[i];
    assert (values[i..n])[1..] == values[i + 1..n];
    assert sum(values[i..n]) == values[i] + sum(values[i + 1..n]);
    total := values[i] + total;
    assert total == sum(values[i..n]);
  }
  assert i == 0;
  assert values[0..n] == values[..];
  assert total == sum(values[..]);

  min := values[0];
  max := values[0];
  var minIndex := 0;
  var maxIndex := 0;
  i := 1;
  while i < n
    invariant 1 <= i <= n
    invariant forall k: int :: 0 <= k < i ==> min <= values[k]
    invariant forall k: int :: 0 <= k < i ==> max >= values[k]
    invariant 0 <= minIndex < i
    invariant min == values[minIndex]
    invariant 0 <= maxIndex < i
    invariant max == values[maxIndex]
  {
    var x := values[i];
    if x < min {
      min := x;
      minIndex := i;
    }
    if x > max {
      max := x;
      maxIndex := i;
    }
    i := i + 1;
  }
  assert i == n;
  assert forall k: int :: 0 <= k < n ==> min <= values[k];
  assert forall k: int :: 0 <= k < n ==> max >= values[k];

  average := total / (n as real);
  range := max - min;
  assert n == values.Length;
  assert average == sum(values[..]) / (values.Length as real);

  forall v: real | v in values[..]
    ensures min <= v
  {
    assert exists k: int :: 0 <= k < |values[..]| && (values[..])[k] == v;
    var k: int :| 0 <= k < |values[..]| && (values[..])[k] == v;
    assert |values[..]| == n;
    assert (values[..])[k] == values[k];
    assert v == values[k];
    assert min <= values[k];
  }
  forall v: real | v in values[..]
    ensures max >= v
  {
    assert exists k: int :: 0 <= k < |values[..]| && (values[..])[k] == v;
    var k: int :| 0 <= k < |values[..]| && (values[..])[k] == v;
    assert |values[..]| == n;
    assert (values[..])[k] == values[k];
    assert v == values[k];
    assert max >= values[k];
  }

  assert 0 <= minIndex < |values[..]|;
  assert (values[..])[minIndex] == min;
  assert min in values[..];
  assert exists v: real :: (v in values[..]) && min == v;

  assert 0 <= maxIndex < |values[..]|;
  assert (values[..])[maxIndex] == max;
  assert max in values[..];
  assert exists v: real :: (v in values[..]) && max == v;
}
\end{lstlisting}
\section*{Total Token Usage}
\textbf{Input tokens: }5878
\\\textbf{Output tokens: }12640
\\\textbf{Reasoning tokens: }9834
\\\textbf{Sum of `total tokens': }18518
\section*{Experiment Timings}
\textbf{Overall Experiment} started at 1780655576606, ended at 1780655867699, lasting 291093ms (291.09 seconds)
\\\textbf{Iteration \#1} started at 1780655576606, ended at 1780655749995, lasting 173389ms (173.39 seconds)
\\\textbf{Iteration \#2} started at 1780655749996, ended at 1780655796534, lasting 46538ms (46.54 seconds)
\\\textbf{Iteration \#3} started at 1780655796534, ended at 1780655867699, lasting 71165ms (71.17 seconds)
\end{document}
